% Options for packages loaded elsewhere
\PassOptionsToPackage{unicode}{hyperref}
\PassOptionsToPackage{hyphens}{url}
\PassOptionsToPackage{dvipsnames,svgnames,x11names}{xcolor}
%
\documentclass[
  10pt,
]{article}
\usepackage{amsmath,amssymb}
\usepackage{iftex}
\ifPDFTeX
  \usepackage[T1]{fontenc}
  \usepackage[utf8]{inputenc}
  \usepackage{textcomp} % provide euro and other symbols
\else % if luatex or xetex
  \usepackage{unicode-math} % this also loads fontspec
  \defaultfontfeatures{Scale=MatchLowercase}
  \defaultfontfeatures[\rmfamily]{Ligatures=TeX,Scale=1}
\fi
\usepackage{lmodern}
\ifPDFTeX\else
  % xetex/luatex font selection
  \setmainfont[]{DejaVu Serif}
  \setsansfont[]{DejaVu Sans}
  \setmonofont[]{DejaVu Sans Mono}
  \setmathfont[]{Latin Modern Math}
\fi
% Use upquote if available, for straight quotes in verbatim environments
\IfFileExists{upquote.sty}{\usepackage{upquote}}{}
\IfFileExists{microtype.sty}{% use microtype if available
  \usepackage[]{microtype}
  \UseMicrotypeSet[protrusion]{basicmath} % disable protrusion for tt fonts
}{}
\makeatletter
\@ifundefined{KOMAClassName}{% if non-KOMA class
  \IfFileExists{parskip.sty}{%
    \usepackage{parskip}
  }{% else
    \setlength{\parindent}{0pt}
    \setlength{\parskip}{6pt plus 2pt minus 1pt}}
}{% if KOMA class
  \KOMAoptions{parskip=half}}
\makeatother
\usepackage{xcolor}
\usepackage[margin=2.0cm]{geometry}
\usepackage{color}
\usepackage{fancyvrb}
\newcommand{\VerbBar}{|}
\newcommand{\VERB}{\Verb[commandchars=\\\{\}]}
\DefineVerbatimEnvironment{Highlighting}{Verbatim}{commandchars=\\\{\}}
% Add ',fontsize=\small' for more characters per line
\newenvironment{Shaded}{}{}
\newcommand{\AlertTok}[1]{\textcolor[rgb]{1.00,0.00,0.00}{\textbf{#1}}}
\newcommand{\AnnotationTok}[1]{\textcolor[rgb]{0.38,0.63,0.69}{\textbf{\textit{#1}}}}
\newcommand{\AttributeTok}[1]{\textcolor[rgb]{0.49,0.56,0.16}{#1}}
\newcommand{\BaseNTok}[1]{\textcolor[rgb]{0.25,0.63,0.44}{#1}}
\newcommand{\BuiltInTok}[1]{\textcolor[rgb]{0.00,0.50,0.00}{#1}}
\newcommand{\CharTok}[1]{\textcolor[rgb]{0.25,0.44,0.63}{#1}}
\newcommand{\CommentTok}[1]{\textcolor[rgb]{0.38,0.63,0.69}{\textit{#1}}}
\newcommand{\CommentVarTok}[1]{\textcolor[rgb]{0.38,0.63,0.69}{\textbf{\textit{#1}}}}
\newcommand{\ConstantTok}[1]{\textcolor[rgb]{0.53,0.00,0.00}{#1}}
\newcommand{\ControlFlowTok}[1]{\textcolor[rgb]{0.00,0.44,0.13}{\textbf{#1}}}
\newcommand{\DataTypeTok}[1]{\textcolor[rgb]{0.56,0.13,0.00}{#1}}
\newcommand{\DecValTok}[1]{\textcolor[rgb]{0.25,0.63,0.44}{#1}}
\newcommand{\DocumentationTok}[1]{\textcolor[rgb]{0.73,0.13,0.13}{\textit{#1}}}
\newcommand{\ErrorTok}[1]{\textcolor[rgb]{1.00,0.00,0.00}{\textbf{#1}}}
\newcommand{\ExtensionTok}[1]{#1}
\newcommand{\FloatTok}[1]{\textcolor[rgb]{0.25,0.63,0.44}{#1}}
\newcommand{\FunctionTok}[1]{\textcolor[rgb]{0.02,0.16,0.49}{#1}}
\newcommand{\ImportTok}[1]{\textcolor[rgb]{0.00,0.50,0.00}{\textbf{#1}}}
\newcommand{\InformationTok}[1]{\textcolor[rgb]{0.38,0.63,0.69}{\textbf{\textit{#1}}}}
\newcommand{\KeywordTok}[1]{\textcolor[rgb]{0.00,0.44,0.13}{\textbf{#1}}}
\newcommand{\NormalTok}[1]{#1}
\newcommand{\OperatorTok}[1]{\textcolor[rgb]{0.40,0.40,0.40}{#1}}
\newcommand{\OtherTok}[1]{\textcolor[rgb]{0.00,0.44,0.13}{#1}}
\newcommand{\PreprocessorTok}[1]{\textcolor[rgb]{0.74,0.48,0.00}{#1}}
\newcommand{\RegionMarkerTok}[1]{#1}
\newcommand{\SpecialCharTok}[1]{\textcolor[rgb]{0.25,0.44,0.63}{#1}}
\newcommand{\SpecialStringTok}[1]{\textcolor[rgb]{0.73,0.40,0.53}{#1}}
\newcommand{\StringTok}[1]{\textcolor[rgb]{0.25,0.44,0.63}{#1}}
\newcommand{\VariableTok}[1]{\textcolor[rgb]{0.10,0.09,0.49}{#1}}
\newcommand{\VerbatimStringTok}[1]{\textcolor[rgb]{0.25,0.44,0.63}{#1}}
\newcommand{\WarningTok}[1]{\textcolor[rgb]{0.38,0.63,0.69}{\textbf{\textit{#1}}}}
\usepackage{longtable,booktabs,array}
\usepackage{calc} % for calculating minipage widths
% Correct order of tables after \paragraph or \subparagraph
\usepackage{etoolbox}
\makeatletter
\patchcmd\longtable{\par}{\if@noskipsec\mbox{}\fi\par}{}{}
\makeatother
% Allow footnotes in longtable head/foot
\IfFileExists{footnotehyper.sty}{\usepackage{footnotehyper}}{\usepackage{footnote}}
\makesavenoteenv{longtable}
\setlength{\emergencystretch}{3em} % prevent overfull lines
\providecommand{\tightlist}{%
  \setlength{\itemsep}{0pt}\setlength{\parskip}{0pt}}
\setcounter{secnumdepth}{-\maxdimen} % remove section numbering
% doc_header.tex -- LaTeX preamble additions used by build_docs.sh for every
% Markdown -> LaTeX -> PDF conversion (supplements, README, submission files).
\usepackage{etoolbox}
% longtable may not be loaded by pandoc when a document has no table
\usepackage{longtable}
% bold vectors: \bm -> \symbf under unicode-math (XeLaTeX), bm package otherwise
\makeatletter
\@ifpackageloaded{unicode-math}{\providecommand{\bm}[1]{\symbf{#1}}}{\usepackage{bm}}
\makeatother
% smaller type in long tables so that wide audit tables fit the text width;
% inside a table a line break is allowed after an underscore of plain text
% (pandoc writes it as \_), so that long identifiers wrap inside their column
% rows are set a little apart, so that symbols with indices do not touch
\AtBeginEnvironment{longtable}{\footnotesize\renewcommand{\arraystretch}{1.45}%
  \let\sptunderscore\_\def\_{\sptunderscore\allowbreak}}
\setlength{\LTcapwidth}{\textwidth}
% allow line breaks after slashes and hyphens inside table cells
\usepackage{xurl}
% code blocks: break over-long lines instead of running into the margin.
% Plain ``` blocks are typeset by pandoc as {verbatim}, highlighted ones
% (```bash) as {Highlighting}; both are re-declared as breakable fancyvrb
% environments.  The continuation line is marked with a small hooked arrow.
\usepackage{fvextra}
\RecustomVerbatimEnvironment{verbatim}{Verbatim}{breaklines,breakanywhere,fontsize=\small}
\makeatletter
\@ifundefined{Highlighting}{}{%
  \RecustomVerbatimEnvironment{Highlighting}{Verbatim}{breaklines,breakanywhere,fontsize=\small,commandchars=\\\{\}}}
\makeatother
% inline code (pandoc writes it as \texttt{...}): never hyphenate -- an inserted
% hyphen could be read as part of a file name -- but allow a line break after
% an underscore, slash, comma, colon or equals sign, so that long paths and
% column lists wrap inside the text width and inside table columns.
% Inside a formula (e.g. the name of a notebook symbol in an equation) the same
% command must first leave math mode, which \text of amsmath does.
\ExplSyntaxOn
\tl_new:N \l__spt_code_tl
\cs_new_protected:Npn \__spt_texttt:n #1
  {
    \tl_set:Nn \l__spt_code_tl {#1}
    \regex_replace_all:nnN { ([\/\,\:\=]) } { \1 \c{allowbreak} } \l__spt_code_tl
    \regex_replace_all:nnN { \c{_} } { \c{_} \c{allowbreak} } \l__spt_code_tl
    \group_begin: \ttfamily \hyphenchar \font = -1 \scan_stop: \l__spt_code_tl \group_end:
  }
\RenewDocumentCommand \texttt { m }
  {
    \mode_if_math:TF { \text { \__spt_texttt:n {#1} } } { \__spt_texttt:n {#1} }
  }
\ExplSyntaxOff
% A figure is written in the Markdown source as an image paragraph followed by
% a caption paragraph ("**Figure S1.** ..."), so that the caption is also
% visible on GitHub.  An image that starts a paragraph is centered.
\usepackage{graphicx}% pandoc loads it only for documents that contain an image
\NewCommandCopy\sptincludegraphics\includegraphics
\RenewDocumentCommand\includegraphics{O{} m}{%
  \ifvmode\begin{center}\sptincludegraphics[#1]{#2}\end{center}%
  \else\sptincludegraphics[#1]{#2}\fi}
% keep figures where they are referenced
\usepackage{float}
\floatplacement{figure}{H}
% running header with the document title
\usepackage{fancyhdr}
\pagestyle{fancy}
\fancyhf{}
\fancyhead[L]{\small PROJECT spacetimes, Revision\_Deflation --- documentation and provenance}
\fancyhead[R]{\small\thepage}
\renewcommand{\headrulewidth}{0.4pt}
% 16x16 matrices: tighter columns so that a full matrix fits the text width
\setlength{\arraycolsep}{3.2pt}
% The documents of this project contain long tagged symbols, file paths and
% identifiers that cannot be hyphenated.  A paragraph is therefore allowed to
% have looser lines before TeX gives up and lets a line stick out of the margin.
\AtBeginDocument{\tolerance=2000 \emergencystretch=5em \hbadness=2000}
\ifLuaTeX
  \usepackage{selnolig}  % disable illegal ligatures
\fi
\IfFileExists{bookmark.sty}{\usepackage{bookmark}}{\usepackage{hyperref}}
\IfFileExists{xurl.sty}{\usepackage{xurl}}{} % add URL line breaks if available
\urlstyle{same}
\hypersetup{
  pdftitle={RP1 --- MASTER provenance of the symbols of Revision\_Deflation},
  pdfauthor={Prepared for Patrick Lee Nash, Ph.D.},
  colorlinks=true,
  linkcolor={blue},
  filecolor={Maroon},
  citecolor={blue},
  urlcolor={blue},
  pdfcreator={LaTeX via pandoc}}

\title{RP1 --- MASTER provenance of the symbols of Revision\_Deflation}
\usepackage{etoolbox}
\makeatletter
\providecommand{\subtitle}[1]{% add subtitle to \maketitle
  \apptocmd{\@title}{\par {\large #1 \par}}{}{}
}
\makeatother
\subtitle{Proof, required by Hard Rule 5, that every symbol of the
revision has one meaning and collides neither with the MASTER registry
of the main project nor with the author's notebooks}
\author{Prepared for Patrick Lee Nash, Ph.D.}
\date{November 12, 2026}

\begin{document}
\maketitle

\hypertarget{what-this-document-proves}{%
\section{1. What this document proves}\label{what-this-document-proves}}

Hard Rule 5 of the task specification: the original files often use the
same symbol for completely different objects, and the project must
resolve this without ambiguity and prove that it has. The main project
does so with a MASTER registry (\texttt{provenance/symbols.csv} in the
folder above; provenance P1) in which every symbol has exactly one
meaning.

The revision introduces new objects --- rates, scale factors, a curved
metric, a clock --- and therefore new symbols. This document is the
extension of the MASTER registry by those symbols. It proves, by six
machine checks whose output is reproduced in Sec. 4:

\begin{enumerate}
\def\labelenumi{\arabic{enumi}.}
\tightlist
\item
  no two entries of the registry of the revision share a name;
\item
  no new name or symbol is a name or symbol of the main registry;
\item
  no new name is a symbol of one of the author's fourteen notebooks
  (1921 symbols);
\item
  in every document of the revision the author's rates are written
  \(\mathrm{g55}\), \(\mathrm{g66}\), \(\mathrm{g77}\) --- upright and
  without subscript --- and every paragraph in which a component of the
  metric \(g_{55}\), \(g_{66}\) or \(g_{77}\) occurs says that it is a
  component of the metric;
\item
  the names that an earlier version of the programs used for new objects
  although the main registry uses them for other objects no longer occur
  in any data file;
\item
  every column of every data file of the revision is a registered name
  or a plain descriptive word.
\end{enumerate}

\hypertarget{the-one-place-where-confusion-was-possible-the-rate-mathrmg55-and-the-component-g_55-of-the-metric}{%
\section{\texorpdfstring{2. The one place where confusion was possible:
the rate \(\mathrm{g55}\) and the component \(g_{55}\) of the
metric}{2. The one place where confusion was possible: the rate \textbackslash mathrm\{g55\} and the component g\_\{55\} of the metric}}\label{the-one-place-where-confusion-was-possible-the-rate-mathrmg55-and-the-component-g_55-of-the-metric}}

The author's names for the three rates are \texttt{g55}, \texttt{g66},
\texttt{g77}. The components of the metric along the same three
coordinates are, in the usual notation, \(g_{55}\), \(g_{66}\),
\(g_{77}\). The two are different objects that occur in the same
formula, \[g_{55}=-\,\mathrm{const5}^2\,e^{-2\,\mathrm{g55}\,x4}.\] The
rule of the revision:

\begin{longtable}[]{@{}
  >{\raggedright\arraybackslash}p{(\columnwidth - 4\tabcolsep) * \real{0.3333}}
  >{\raggedright\arraybackslash}p{(\columnwidth - 4\tabcolsep) * \real{0.3333}}
  >{\raggedright\arraybackslash}p{(\columnwidth - 4\tabcolsep) * \real{0.3333}}@{}}
\toprule\noalign{}
\begin{minipage}[b]{\linewidth}\raggedright
object
\end{minipage} & \begin{minipage}[b]{\linewidth}\raggedright
in documents
\end{minipage} & \begin{minipage}[b]{\linewidth}\raggedright
in programs and data files
\end{minipage} \\
\midrule\noalign{}
\endhead
\bottomrule\noalign{}
\endlastfoot
the rate of deflation of \(x5\) & \(\mathrm{g55}\) (upright, no
subscript) & \texttt{g55} \\
the component of the metric & \(g_{55}\) (italic, with subscript),
always called a component of the metric in the same paragraph &
\texttt{g\_55}, or the scale factor \texttt{b\_5} with
\(g_{55}=-b_5^2\) \\
\end{longtable}

and likewise for \(x6\) and \(x7\). In a statement that holds for all
three coordinates the rate is written \(\mathrm{rate}_t\) and the scale
factor \(b_t\), \(t=5,6,7\). Check RS4 enforces the rule in every
document.

\emph{What the notebooks use.} The names \texttt{g55}, \texttt{g66},
\texttt{g77} occur in none of the fourteen notebooks. The notebooks use
\texttt{g44} and \texttt{g71} as names of whole metrics (which the main
project renamed, error E7 of its error log), and those of the notebooks
that contain exponential scale factors write the rate as \texttt{H}. No
collision arises.

\hypertarget{the-registry-of-the-revision}{%
\section{3. The registry of the
revision}\label{the-registry-of-the-revision}}

The registry is the file \texttt{provenance/symbols\_deflation.csv},
written by \texttt{code/scripts/write\_symbol\_registry.py}. ``Origin''
says whether a symbol is the author's, is taken over from the main
project, or is new. ``Defined in'' names the document and the number of
its section: DMD is the document ``dirac matrices deflating'', RS1 to
RS4 are the supplements, RP0 is the provenance of the hypotheses, and M
stands for the two manuscripts, whose sections carry the same numbers in
both (M II is Sec. II of either manuscript; ``(4+4)'' or ``(7+1)'' marks
a symbol that only one of them uses).

\emph{Index letters.} An index letter has no meaning of its own; its
range is stated where it is used. The revision keeps the rule of the
main registry and adds one letter. The indices \(a,b\) label the four
extra axes \(x0,x5,x6,x7\) in the supplements RS1 to RS4 and in the
manuscripts, and the four flat coordinates \(x0,x1,x2,x3\) in DMD and in
the statements of the checks, as in the corresponding documents of the
main project. The two uses never meet in one formula. In the
manuscripts, where both sets of coordinates occur, the flat coordinates
carry the index \(j\) instead, with the range \(j\in\{0,1,2,3\}\) stated
in their introduction. The registry lists the index letters in its first
group.

\textbf{D1 Hypotheses and geometry}

\begin{longtable}[]{@{}
  >{\raggedright\arraybackslash}p{(\columnwidth - 8\tabcolsep) * \real{0.1379}}
  >{\raggedright\arraybackslash}p{(\columnwidth - 8\tabcolsep) * \real{0.1379}}
  >{\raggedright\arraybackslash}p{(\columnwidth - 8\tabcolsep) * \real{0.4483}}
  >{\raggedright\arraybackslash}p{(\columnwidth - 8\tabcolsep) * \real{0.1379}}
  >{\raggedright\arraybackslash}p{(\columnwidth - 8\tabcolsep) * \real{0.1379}}@{}}
\toprule\noalign{}
\begin{minipage}[b]{\linewidth}\raggedright
symbol
\end{minipage} & \begin{minipage}[b]{\linewidth}\raggedright
in programs and data files
\end{minipage} & \begin{minipage}[b]{\linewidth}\raggedright
meaning
\end{minipage} & \begin{minipage}[b]{\linewidth}\raggedright
defined in
\end{minipage} & \begin{minipage}[b]{\linewidth}\raggedright
origin
\end{minipage} \\
\midrule\noalign{}
\endhead
\bottomrule\noalign{}
\endlastfoot
\(\mathrm{g55}\), \(\mathrm{g66}\), \(\mathrm{g77}\) & \texttt{g55},
\texttt{g66}, \texttt{g77} & the three rates of deflation of \(x5\),
\(x6\), \(x7\) (positive constants, dimension 1/length). These are the
author's names; they are written upright and never with a subscript.
They are \emph{not} components of the metric: the components are
\(g_{55}\), \(g_{66}\), \(g_{77}\) & RP0 1; DMD 0.1; M I & author \\
\(\mathrm{rate}_t\) & \texttt{rate\_t}, \texttt{rate} & the rate of the
coordinate \(xt\) in a statement that holds for \(t=5,6,7\) alike:
\(\mathrm{rate}_5=\mathrm{g55}\), \(\mathrm{rate}_6=\mathrm{g66}\),
\(\mathrm{rate}_7=\mathrm{g77}\) & DMD 0.1; M I & new \\
\(\mathrm{gsum}\) & \texttt{gsum}, \texttt{gam} &
\(\mathrm{g55}+\mathrm{g66}+\mathrm{g77}\), the total rate: the proper
volume of the three deflating coordinates falls as
\(e^{-\mathrm{gsum}\,x4}\) (\texttt{gam}: gsum in units of 100 km/s/Mpc)
& DMD 1.1; M II & new \\
\(\mathrm{gsq}\), \(\mathrm{gprod}\) & \texttt{gsq}, \texttt{gprod} &
\(\mathrm{g55}^2+\mathrm{g66}^2+\mathrm{g77}^2\) and
\(\mathrm{g55}\,\mathrm{g66}+\mathrm{g55}\,\mathrm{g77}+\mathrm{g66}\,\mathrm{g77}\);
\(\mathrm{gsum}^2=\mathrm{gsq}+2\,\mathrm{gprod}\) & DMD 1.3; M II &
new \\
\(\mathrm{const5}\), \(\mathrm{const6}\), \(\mathrm{const7}\);
\(\mathrm{const}_t\) & \texttt{const5}, \texttt{const6},
\texttt{const7}, \texttt{const\_t}, \texttt{cst} & the three positive
constants of the author's metric (his word: const): the values of the
scale factors at \(x4=0\) & RP0 1; DMD 0.1; M I & author \\
\(b_5\), \(b_6\), \(b_7\); \(b_t\) & \texttt{b\_5}, \texttt{b\_6},
\texttt{b\_7}, \texttt{b\_t} & the scale factors of the deflating
coordinates, \(b_t(x4)=\mathrm{const}_t\,e^{-\mathrm{rate}_t\,x4}\): a
coordinate distance \(dxt\) is the proper length
\(b_t\,\lvert dxt\rvert\). Always written with the index 5, 6, 7 or
\(t\) (the letter \(b\) without index is a Ginzburg-Landau coefficient
of the main registry) & DMD 1.1; M II & new \\
\(V_d\) & \texttt{Vd} & \(b_5b_6b_7=\sqrt{\lvert\det g\rvert}\), the
proper volume of a unit coordinate cell of the deflating coordinates;
\(V_d\propto e^{-\mathrm{gsum}\,x4}\) & DMD 1.1; RS1 1; M II & new \\
\(g_{1644d}\), \(g_{1671d}\); \(g_{\mu\nu}\) & \texttt{g\_1644d},
\texttt{g\_1671d} & the metrics of the deflating 4+4 and 7+1 spacetimes,
\(\mathrm{diag}(\eta_{00},\dots,\eta_{44},\eta_{55}b_5^2,\eta_{66}b_6^2,\eta_{77}b_7^2)\)
with \(\eta=\eta_{1644}\) or \(\eta_{1671}\); in both,
\(g_{tt}=-\mathrm{const}_t^2e^{-2\,\mathrm{rate}_t\,x4}\) & DMD 1.1; M
II & author \\
\(\mu,\nu,\rho,\sigma\) (indices) & \texttt{mu}, \texttt{nu},
\texttt{rho}, \texttt{sigma}, \texttt{lam} & coordinate indices with the
values \(0,\dots,7\) on the deflating spacetimes (only as sub- or
superscripts; as stand-alone letters \(\mu\), \(\nu\), \(\lambda\) keep
the meanings of the main registry). The capital indices \(A,B,C,D\) of
the main registry, also with a prime, label the orthonormal frame: the
frame vector \(A\) points along the coordinate \(xA\) (as a stand-alone
letter \(C\) is the coupling of the main registry) & DMD 1.1; M II &
new \\
\(t\), \(s\) (indices) & --- & indices of the three deflating
coordinates, values 5, 6, 7: \(t\) as in the main registry, and \(s\) as
a second such index where two are needed (\(s_h\) with the index \(h\)
are matrices of the main registry) & DMD 0.1, 1.3; M I, II & main \\
\(a,b\); \(j\) (indices) & --- & \(a,b\) as in the main registry, with
the range stated at each place: in RS1 to RS4 and in the manuscripts the
index of the four extra axes, values 0, 5, 6, 7 (on the deflating
spacetimes a frame index); in DMD and in the statements of the checks of
the programs the index of the four flat coordinates other than the time,
values 0, 1, 2, 3. In the manuscripts of the revision, where both sets
of coordinates occur, the flat coordinates other than the time carry the
index \(j\) with the range \(0,1,2,3\) stated there; a sum over \(j\)
with other limits (the ordinary space \(1,2,3\) of the main registry; in
DMD Eq. (3.12) the seven space-like directions) states them & DMD 0.1;
RS1 5.1; M I & main \\
\(e_A{}^\mu\), \(e^A{}_\mu\) & \texttt{e\_A\^{}mu} & the orthonormal
frame and its inverse:
\(e^A{}_\mu=\mathrm{diag}(1,1,1,1,1,b_5,b_6,b_7)\) & DMD 1.1; M II &
new \\
\(\Gamma^\rho{}_{\mu\nu}\) & \texttt{Gamma}, \texttt{christoffel} & the
Christoffel symbols of the deflating metric (three indices; the letter
\(\Gamma\) with square brackets is the effective action of the main
registry) & DMD 1.2; M II & new \\
\(R^\rho{}_{\sigma\mu\nu}\), \(R^\mu{}_\nu\), \(R\) & \texttt{riemann},
\texttt{ricci}, \texttt{Ric44}, \texttt{Ric55}, \texttt{Ric66},
\texttt{Ric77}, \texttt{Rscalar} & Riemann tensor, Ricci tensor and
curvature scalar of the deflating metric & DMD 1.3; M II & new \\
\(\mathcal G^\mu{}_\nu\) & \texttt{Ein00}, \texttt{Ein44},
\texttt{Ein55}, \texttt{Ein66}, \texttt{Ein77}, \texttt{einstein\_mixed}
& the Einstein tensor \(R^\mu{}_\nu-\tfrac12R\,\delta^\mu_\nu\) (script
letter, so that it cannot be taken for Newton's constant or for a
constant of the Fermi type) & DMD 1.4; M II & new \\
\(\mathcal X\); \(\mathcal X_{\rm scale}\), \(\mathcal X_{(s)}\) &
\texttt{killing\_vectors}, \texttt{killing\_count} & a Killing vector
field of the deflating metric (script letter; the plain letter \(X\) is
the set of coordinates of the main registry); the scaling vector
\(\partial_4+\sum_t\mathrm{rate}_t\,xt\,\partial_t\); the three special
Killing vectors of the case of equal rates, \(s=5,6,7\) & DMD 1.5, 1.6;
M II, App. B & new \\
\(\bar x4\); \(\bar x5\), \(\bar x6\), \(\bar x7\) & --- & coordinates
of the case of equal rates in which the line element of \(x4,\dots,x7\)
is a multiple of a flat one:
\(\bar x4=e^{\mathrm{g55}\,x4}/\mathrm{g55}\),
\(\bar xt=\mathrm{const}_t\,xt\) & DMD 1.6; M App. B & new \\
\(\mathcal H^\mu{}_\nu\); \(\alpha_{GB}\), \(\Lambda_8\) &
\texttt{Lanczos}, \texttt{alpha\_GB}, \texttt{Lambda\_8} & the Lanczos
tensor of the Gauss-Bonnet term (script letter; \(H\) is the Hubble
rate); the coupling of the Gauss-Bonnet term and the cosmological
constant of Einstein-Gauss-Bonnet gravity in eight dimensions (with the
index 8; \(\Lambda\) without index is the momentum cutoff of the main
registry) & DMD 1.4a; M II & new \\
\(K_{AB}\) & \texttt{ksec} & the sectional curvature of the plane of the
frame directions \(A\) and \(B\), in the convention in which the
space-like components of the metric are positive:
\(K_{4t}=\mathrm{rate}_t^2\),
\(K_{ts}=\mathrm{rate}_t\,\mathrm{rate}_s\) for \(t\neq s\), zero if one
direction is flat. Always with two subscripts: the letter \(K\) without
subscript is the generator of the main registry, and \(K\) with an order
as subscript and with an argument is a modified Bessel function & DMD
1.4a; M App. B & new \\
\(\varsigma\) & \texttt{affine} & affine parameter along a geodesic; for
a massive particle on the 7+1 spacetime, the proper time divided by the
mass & DMD 1.7; M II & new \\
\(P_\mu\) & \texttt{P\_1}, \texttt{P\_5} & the conserved covariant
momenta of a geodesic (with an index; \(P\) without index is the
generator of the main registry) & DMD 1.7; M II & new \\
\(x4_{\rm turn}\) & \texttt{x4\_turn} & the value of the ordinary time
at which a classical particle with momentum along a deflating time turns
back, and beyond which a mode that depends on a deflating time is
extinguished (4+4): there \(k_d=E_\perp\) & DMD 1.7, 4.3; M II (4+4) &
new \\
\end{longtable}

\textbf{D2 Matrices, fields and their equations}

\begin{longtable}[]{@{}
  >{\raggedright\arraybackslash}p{(\columnwidth - 8\tabcolsep) * \real{0.1379}}
  >{\raggedright\arraybackslash}p{(\columnwidth - 8\tabcolsep) * \real{0.1379}}
  >{\raggedright\arraybackslash}p{(\columnwidth - 8\tabcolsep) * \real{0.4483}}
  >{\raggedright\arraybackslash}p{(\columnwidth - 8\tabcolsep) * \real{0.1379}}
  >{\raggedright\arraybackslash}p{(\columnwidth - 8\tabcolsep) * \real{0.1379}}@{}}
\toprule\noalign{}
\begin{minipage}[b]{\linewidth}\raggedright
symbol
\end{minipage} & \begin{minipage}[b]{\linewidth}\raggedright
in programs and data files
\end{minipage} & \begin{minipage}[b]{\linewidth}\raggedright
meaning
\end{minipage} & \begin{minipage}[b]{\linewidth}\raggedright
defined in
\end{minipage} & \begin{minipage}[b]{\linewidth}\raggedright
origin
\end{minipage} \\
\midrule\noalign{}
\endhead
\bottomrule\noalign{}
\endlastfoot
\(\gamma_{1644d}^\mu\), \(\gamma_{1671d}^\mu\) & \texttt{gamma\_1644d},
\texttt{gamma\_1671d} & the Dirac matrices of the deflating spacetimes
in coordinate components, \(\gamma^\mu=e_A{}^\mu\gamma^A\):
\(\gamma^t=\gamma^{A=t}/b_t\) for \(t=5,6,7\), the other five are the
constant matrices \(\gamma_{1644}^A\), \(\gamma_{1671}^A\) of the main
project & DMD 2.2; M III & new \\
\(\omega_{\mu AB}\) & \texttt{omega\_(mu}, \texttt{A}, \texttt{B)} & the
coefficients of the spin connection (three indices; \(\omega_{1644}\),
\(\omega_{1671}\) with a tag are the volume elements of the main
registry) & DMD 2.3; M III & new \\
\(\Omega_\mu\) & \texttt{Omega\_mu}, \texttt{omega} & the spin
connection as a matrix, \(\tfrac14\omega_{\mu AB}\gamma^A\gamma^B\); the
covariant derivative of a spinor is
\(\nabla_\mu\psi=(\partial_\mu+\Omega_\mu)\psi\) & DMD 2.3; M III &
new \\
\(\psi_{1644d}\), \(\psi_{1671d}\) & \texttt{dirac1644d},
\texttt{dirac1671d} & the fields dirac1644d and dirac1671d: the
16-component fields of the main project on the deflating spacetimes;
\((\gamma^\mu\nabla_\mu-m)\psi_{1644d}=0\) and
\((i\gamma^\mu\nabla_\mu-m)\psi_{1671d}=0\) & DMD 3.1, 3.2; M III &
new \\
\(\hat\psi\) & \texttt{psihat} & the weighted field
\(\sqrt{V_d}\,\psi\), which obeys the equation without spin connection &
DMD 3.3; M III & new \\
\(K(x4)\), \(\mathsf h(x4)\) & \texttt{K}, \texttt{h} & the generator
\(K_{1644}\) and the Hamiltonian matrix \(\mathsf h_{1671}\) of the main
registry with the wavenumbers along the deflating coordinates divided by
the scale factors: \(\partial_4\hat\psi=K\hat\psi\) (4+4),
\(i\partial_4\hat\psi=\mathsf h\hat\psi\) (7+1) & DMD 3.4; M III &
main \\
\end{longtable}

\textbf{D3 Solutions mode by mode}

\begin{longtable}[]{@{}
  >{\raggedright\arraybackslash}p{(\columnwidth - 8\tabcolsep) * \real{0.1379}}
  >{\raggedright\arraybackslash}p{(\columnwidth - 8\tabcolsep) * \real{0.1379}}
  >{\raggedright\arraybackslash}p{(\columnwidth - 8\tabcolsep) * \real{0.4483}}
  >{\raggedright\arraybackslash}p{(\columnwidth - 8\tabcolsep) * \real{0.1379}}
  >{\raggedright\arraybackslash}p{(\columnwidth - 8\tabcolsep) * \real{0.1379}}@{}}
\toprule\noalign{}
\begin{minipage}[b]{\linewidth}\raggedright
symbol
\end{minipage} & \begin{minipage}[b]{\linewidth}\raggedright
in programs and data files
\end{minipage} & \begin{minipage}[b]{\linewidth}\raggedright
meaning
\end{minipage} & \begin{minipage}[b]{\linewidth}\raggedright
defined in
\end{minipage} & \begin{minipage}[b]{\linewidth}\raggedright
origin
\end{minipage} \\
\midrule\noalign{}
\endhead
\bottomrule\noalign{}
\endlastfoot
\(k_5\), \(k_6\), \(k_7\); \(k_t\) & \texttt{k\_5}, \texttt{k\_6},
\texttt{k\_7}, \texttt{k\_t} & the conserved wavenumbers of a mode along
the deflating coordinates (covariant components of the wave vector
\(k_A\) of the main project) & DMD 3.4; M III & new \\
\(E_\perp\) & \texttt{E\_perp} & the energy of a mode, or of a classical
particle, without its motion along the deflating coordinates:
\(\sqrt{m^2+k_0^2+k_1^2+k_2^2+k_3^2}\) (with the wavenumbers or momenta
along \(x0,\dots,x3\)) & DMD 1.7, 3.4; M II, III & new \\
\(k_d\) & \texttt{k\_d} & the physical wavenumber along the deflating
coordinates, \(k_d(x4)=\big[\sum_tk_t^2/b_t(x4)^2\big]^{1/2}\); it grows
as the coordinates deflate & DMD 3.4; M III & new \\
\(\mathsf h_\perp\), \(K_\perp\) & \texttt{h\_perp}, \texttt{K\_perp} &
the parts of \(\mathsf h\) and of \(K\) that do not contain the scale
factors: the operators of the modes that do not depend on the deflating
coordinates, \(\mathsf h_\perp=m\beta+\sum_ap_a\alpha^a\) with
\(\mathsf h_\perp^2=E_\perp^2\), and
\(K_\perp=-m\gamma^4+i\sum_ak_a\gamma^4\gamma^a\) with
\(K_\perp^2=-E_\perp^2\) & DMD 4.1; M III & new \\
\(\gamma^d\), \(\alpha^d\) & --- & the frame matrix along the direction
of the wave vector in the deflating coordinates,
\(\gamma^d=k_d^{-1}\sum_t(k_t/b_t)\gamma^t\), and
\(\alpha^d=\gamma^4\gamma^d\); constant in time when one rate governs &
DMD 4.1; M III & new \\
\(w\); \(f_+\), \(f_-\); \(Z_+\), \(Z_-\) & --- & auxiliary objects of
the exact solutions (DMD 4.2 and 4.3; \(w\), \(Z_+\), \(Z_-\) also in
Sec. IV of the manuscripts): a constant eigenvector of
\(\mathsf h_\perp\) or of \(K_\perp\); the two coefficient functions of
the ansatz; two cylinder functions whose orders differ by one & DMD 4.2,
4.3; M IV & new \\
\(u_d\) & \texttt{u\_d} & \(k_d/\mathrm{rate}_t\) for a mode that
depends on one deflating coordinate: the argument of the Bessel
functions of the exact solutions & DMD 4.1; M IV & new \\
\(n_{\rm cr}\) & \texttt{n\_cr}, \texttt{n\_min}, \texttt{n\_mean},
\texttt{n\_max}, \texttt{n\_exact}, \texttt{fd5}, \texttt{fd6} & the
number of particles per internal state that the deflation creates in a
mode that starts empty (7+1) & DMD 4.2; M IV (7+1) & new \\
\(T_{\rm defl}\) & --- & the temperature of the thermal spectrum of the
created quanta, \(\mathrm{rate}_t/2\pi\) in units with \(\hbar=k_B=1\)
(7+1) & DMD 4.2; M IV (7+1) & new \\
\(S_{\rm WKB}\) & \texttt{S\_wkb} & \(\int\sqrt{k_d^2-E_\perp^2}\,dx4\)
from the turning value on: exponent of the extinction of a bounded
solution (4+4) & DMD 4.3; M IV (4+4) & new \\
\(\varphi_{\rm bos}\), \(\hat\varphi_{\rm bos}\), \(m_{\rm bos}\) &
\texttt{phi\_bos}, \texttt{Shat}, \texttt{Shat2}, \texttt{m\_bos} & a
mode of the Vorton-Scalar (the radial mode, \(m_{\rm bos}=m_s\), or the
phase mode, \(m_{\rm bos}=0\)), its weighted form
\(\sqrt{V_d}\,\varphi_{\rm bos}\), and its mass & RS1 2; M V & new \\
\(\varpi_b\) & \texttt{varpi\_b} &
\(\sqrt{k^2+m_{\rm bos}^2-\mathrm{gsum}^2/4}\): the frequency of the
weighted boson mode that does not depend on the deflating coordinates &
RS1 2; M V & new \\
\end{longtable}

\textbf{D4 The deflation clock and the couplings}

\begin{longtable}[]{@{}
  >{\raggedright\arraybackslash}p{(\columnwidth - 8\tabcolsep) * \real{0.1379}}
  >{\raggedright\arraybackslash}p{(\columnwidth - 8\tabcolsep) * \real{0.1379}}
  >{\raggedright\arraybackslash}p{(\columnwidth - 8\tabcolsep) * \real{0.4483}}
  >{\raggedright\arraybackslash}p{(\columnwidth - 8\tabcolsep) * \real{0.1379}}
  >{\raggedright\arraybackslash}p{(\columnwidth - 8\tabcolsep) * \real{0.1379}}@{}}
\toprule\noalign{}
\begin{minipage}[b]{\linewidth}\raggedright
symbol
\end{minipage} & \begin{minipage}[b]{\linewidth}\raggedright
in programs and data files
\end{minipage} & \begin{minipage}[b]{\linewidth}\raggedright
meaning
\end{minipage} & \begin{minipage}[b]{\linewidth}\raggedright
defined in
\end{minipage} & \begin{minipage}[b]{\linewidth}\raggedright
origin
\end{minipage} \\
\midrule\noalign{}
\endhead
\bottomrule\noalign{}
\endlastfoot
\(\tau_d\) & \texttt{tau\_d} & the deflation clock,
\(\tau_d=\mathrm{gsum}\,(x4-x4_{\rm ref})\): the number of e-foldings by
which the volume \(V_d\) has shrunk since the reference time (the letter
\(\tau\) without the index \(d\) is the optical depth of the main
registry) & RS1 1; M V & new \\
\(\tau_{\rm pair}\) & \texttt{tau\_pair} & the clock measured from the
moment at which two fermions become bound: \(\ln(D_{VS}/D_c)\) & RS1 4;
RS4 1; M VII & new \\
\(\hat\Phi_{VS}\) & \texttt{Phihat} & the weighted Vorton-Scalar
\(\sqrt{V_d}\,\Phi_{VS}\) & RS1 1; M V & new \\
\end{longtable}

\textbf{D5 Dynamics of the paired state}

\begin{longtable}[]{@{}
  >{\raggedright\arraybackslash}p{(\columnwidth - 8\tabcolsep) * \real{0.1379}}
  >{\raggedright\arraybackslash}p{(\columnwidth - 8\tabcolsep) * \real{0.1379}}
  >{\raggedright\arraybackslash}p{(\columnwidth - 8\tabcolsep) * \real{0.4483}}
  >{\raggedright\arraybackslash}p{(\columnwidth - 8\tabcolsep) * \real{0.1379}}
  >{\raggedright\arraybackslash}p{(\columnwidth - 8\tabcolsep) * \real{0.1379}}@{}}
\toprule\noalign{}
\begin{minipage}[b]{\linewidth}\raggedright
symbol
\end{minipage} & \begin{minipage}[b]{\linewidth}\raggedright
in programs and data files
\end{minipage} & \begin{minipage}[b]{\linewidth}\raggedright
meaning
\end{minipage} & \begin{minipage}[b]{\linewidth}\raggedright
defined in
\end{minipage} & \begin{minipage}[b]{\linewidth}\raggedright
origin
\end{minipage} \\
\midrule\noalign{}
\endhead
\bottomrule\noalign{}
\endlastfoot
\(\Xi_d\) & \texttt{Xi\_d} &
\(\mathrm{gsum}/(\lambda\,\Delta_{\rm eq})\): decides whether the paired
state can follow the deflation (\(\Xi_d\ll1\): yes) & RS1 7; M VIII &
new \\
\(n_k\) & --- & occupation of the pair of states of the single-particle
energy \(\xi_k\) in the shell model of the paired state (with the index
\(k\); the letter \(n\) without index is a density) & RS1 7; M VIII &
new \\
\(\Delta_{\rm eq}\) & \texttt{gap\_eq} & the equilibrium gap for the
pairing strength of the moment & RS1 7; M VIII & new \\
\(E_{\rm sh}\) & \texttt{E\_sh} & half-width in energy of the shell of
states around the Fermi surface in the model of the pair dynamics; unit
of energy and of inverse time there & RS1 7; M VIII & new \\
\(n_{\rm ex}\) & \texttt{n\_ex}, \texttt{n\_ex\_final} & density of
broken pairs (in units of \(E_{\rm sh}\) times the density of states) &
RS1 7; M VIII & new \\
\end{longtable}

\textbf{D6 Cosmology and bounds}

\begin{longtable}[]{@{}
  >{\raggedright\arraybackslash}p{(\columnwidth - 8\tabcolsep) * \real{0.1379}}
  >{\raggedright\arraybackslash}p{(\columnwidth - 8\tabcolsep) * \real{0.1379}}
  >{\raggedright\arraybackslash}p{(\columnwidth - 8\tabcolsep) * \real{0.4483}}
  >{\raggedright\arraybackslash}p{(\columnwidth - 8\tabcolsep) * \real{0.1379}}
  >{\raggedright\arraybackslash}p{(\columnwidth - 8\tabcolsep) * \real{0.1379}}@{}}
\toprule\noalign{}
\begin{minipage}[b]{\linewidth}\raggedright
symbol
\end{minipage} & \begin{minipage}[b]{\linewidth}\raggedright
in programs and data files
\end{minipage} & \begin{minipage}[b]{\linewidth}\raggedright
meaning
\end{minipage} & \begin{minipage}[b]{\linewidth}\raggedright
defined in
\end{minipage} & \begin{minipage}[b]{\linewidth}\raggedright
origin
\end{minipage} \\
\midrule\noalign{}
\endhead
\bottomrule\noalign{}
\endlastfoot
\(G_N\), \(G_8\) & \texttt{G\_N}, \texttt{GN\_early}, \texttt{GN\_star},
\texttt{G\_8} & Newton's constant of ordinary space, and the
gravitational coupling of the eight-dimensional Einstein theory
(assumption A2):
\(G_N=G_8/(\text{proper volume of the extra coordinates})\propto e^{\mathrm{gsum}\,x4}\)
& RS3 1; M XI & new \\
\(a_{\rm sp}\), \(H\) & --- & scale factor of ordinary space (assumption
A1) and its rate \(H=(da_{\rm sp}/dx4)/a_{\rm sp}\); \(H_0\), \(h\) as
in the main registry. (The author's Einstein-Lovelock notebooks use the
letter \(H\) for the rate in the exponent of their metric; that rate
corresponds here to \(\mathrm{g55}\), \(\mathrm{g66}\),
\(\mathrm{g77}\)) & RS3 1; M XI & new \\
\(t_0\) & \texttt{age}, \texttt{age\_Gyr} & the age of the universe in
the model at hand & RS3 2; M XI & new \\
\(\theta_*\), \(z_*\) & \texttt{theta\_star}, \texttt{z\_star} & the
angle under which the sound horizon at last scattering is seen,
\(r_s(z_*)/D_A(z_*)\) with the comoving distance \(D_A\), and the
redshift \(z_*\) of last scattering at which the optical depth without
reionization is 1 (the \(100\,\theta_*\) of the Planck papers). The
angle \(\theta_s\) of the main registry is the same ratio at the maximum
of the visibility function; for the model without deflation
\(100\,\theta_*=1.03965\) and \(100\,\theta_s=1.040358\) & RS3 1, 2; M
XI & new \\
\(\delta_m\), \(f_g\), \(\sigma_8\) & \texttt{growth},
\texttt{f\_growth}, \texttt{growth\_ratio}, \texttt{f\_growth\_ratio},
\texttt{f\_sigma8\_ratio} & density contrast of matter in ordinary
space; its growth rate \(d\ln\delta_m/d\ln a_{\rm sp}\) today; the
amplitude of clustering (only ratios of \(\sigma_8\) are used) & RS3 5;
M XI & new \\
\(v_{\rm EW}\) (main project), \(n_{\rm EW}\) & \texttt{v\_EW},
\texttt{n\_EW} & the electroweak scale of the supplement S2 of the main
project, and the exponent with which it would follow the coupling under
the hypothesis H2: \(n_{\rm EW}=d\ln v_{\rm EW}/d\ln\alpha_{VS}\) & RS2
5; M XII & new \\
\(\varrho_c\) & \texttt{rho\_c} & abbreviation in the formula for
\(n_{\rm EW}\) of the scale-invariant ladder equation:
\(\varrho_c=\tfrac12\sqrt{\alpha_{VS}/\alpha_c-1}\) (the letter \(\rho\)
without index is a density) & RS2 5; M XII & new \\
\(m_p/m_e\) & \texttt{mu\_pe} & the ratio of the proton mass to the
electron mass (always written as this ratio in the documents;
\texttt{mu\_pe} in the statements of the programs) & RS2 5; RS3 3; M XI
& new \\
\(G_{\rm Fermi}\), \(\alpha_{\rm em}\) & \texttt{G\_Fermi} & the Fermi
constant of the weak interaction and the fine-structure constant (they
enter only in the discussion of bounds) & RS3 1, 3; M XI & new \\
\end{longtable}

\textbf{D7 Symmetries and vortex lines}

\begin{longtable}[]{@{}
  >{\raggedright\arraybackslash}p{(\columnwidth - 8\tabcolsep) * \real{0.1379}}
  >{\raggedright\arraybackslash}p{(\columnwidth - 8\tabcolsep) * \real{0.1379}}
  >{\raggedright\arraybackslash}p{(\columnwidth - 8\tabcolsep) * \real{0.4483}}
  >{\raggedright\arraybackslash}p{(\columnwidth - 8\tabcolsep) * \real{0.1379}}
  >{\raggedright\arraybackslash}p{(\columnwidth - 8\tabcolsep) * \real{0.1379}}@{}}
\toprule\noalign{}
\begin{minipage}[b]{\linewidth}\raggedright
symbol
\end{minipage} & \begin{minipage}[b]{\linewidth}\raggedright
in programs and data files
\end{minipage} & \begin{minipage}[b]{\linewidth}\raggedright
meaning
\end{minipage} & \begin{minipage}[b]{\linewidth}\raggedright
defined in
\end{minipage} & \begin{minipage}[b]{\linewidth}\raggedright
origin
\end{minipage} \\
\midrule\noalign{}
\endhead
\bottomrule\noalign{}
\endlastfoot
\(P_c\) & \texttt{Pc} & \(i\alpha^1\alpha^2\alpha^3\chi\), the product
of the five matrices of the sector that does not depend on the extra
coordinates; equal, up to sign, to the generator \(P\) of \(U(1)_E\) of
the main registry & RS1 6; M X & new \\
\(\theta_{12}\) & --- & the angle around a vortex line that lies along
\(x3\), in the plane of \(x1\) and \(x2\) (with the index 12; \(\theta\)
with other indices and the scattering angle keep the meanings of the
main project) & RS1 6; M X & new \\
\(\mathsf M\) (main registry) & \texttt{M\_gap} & the \(8\times8\) gap
matrix of the main registry; \texttt{M\_gap} is its spelling in the
statements of the programs of the revision & RS1 5; M X & main \\
\(w_z\) & \texttt{w\_z} & a constant spinor of a fermion zero mode on a
vortex line & RS1 6; M X & new \\
\end{longtable}

\textbf{D8 Tags and identifiers}

\begin{longtable}[]{@{}
  >{\raggedright\arraybackslash}p{(\columnwidth - 8\tabcolsep) * \real{0.1379}}
  >{\raggedright\arraybackslash}p{(\columnwidth - 8\tabcolsep) * \real{0.1379}}
  >{\raggedright\arraybackslash}p{(\columnwidth - 8\tabcolsep) * \real{0.4483}}
  >{\raggedright\arraybackslash}p{(\columnwidth - 8\tabcolsep) * \real{0.1379}}
  >{\raggedright\arraybackslash}p{(\columnwidth - 8\tabcolsep) * \real{0.1379}}@{}}
\toprule\noalign{}
\begin{minipage}[b]{\linewidth}\raggedright
symbol
\end{minipage} & \begin{minipage}[b]{\linewidth}\raggedright
in programs and data files
\end{minipage} & \begin{minipage}[b]{\linewidth}\raggedright
meaning
\end{minipage} & \begin{minipage}[b]{\linewidth}\raggedright
defined in
\end{minipage} & \begin{minipage}[b]{\linewidth}\raggedright
origin
\end{minipage} \\
\midrule\noalign{}
\endhead
\bottomrule\noalign{}
\endlastfoot
1644d, 1671d & \texttt{1644d}, \texttt{1671d} & tags of the deflating
4+4 and 7+1 spacetimes and of every object defined on them (the main
project's tags 1644, 1671 with the letter d) & DMD 0.1; M I & new \\
\texttt{44d-..}, \texttt{71d-..}, \texttt{def-..} & --- & identifiers of
the machine checks of the revision: on the deflating 4+4 spacetime, on
the deflating 7+1 spacetime, and for results that hold on both
(\texttt{44-..}, \texttt{71-..} without d are the checks of the main
project, executed again) & DMD Appendix A & new \\
RP0 to RP9, RS1 to RS4, DMD & --- & the provenances and supplements of
the revision and the document "dirac matrices deflating" (P0 to P11, S1
to S4, DM are the documents of the main project) & README & new \\
A1, A2, H2, HV & --- & A1, A2: the two assumptions of the cosmological
part (ordinary space expands; gravity is the eight-dimensional Einstein
theory). H2: the author's hypothesis of the main project that the
superfluid state is the Higgs field. HV: the different assumption that
the vacuum value of the Vorton-Scalar sets the electroweak scale & RP0
3; RS2 1; M XI, XII & new \\
\end{longtable}

\textbf{Letters that the revision deliberately does not use}, because
the main registry or the notebooks use them for other objects:
\(\varepsilon\) and \(\epsilon\) (an energy density; a transformation
parameter), \(\nu\) and \(N\) as free symbols (the number of internal
states; a density), \(C\) (a coupling), \(z\) (a redshift), \(\xi\) (a
single-particle energy), \(\sigma\) without index (a self-energy),
\(\tau\) without index (the optical depth), \(\Lambda\) without index (a
cutoff), and the names \texttt{h0}, \texttt{E0}, \texttt{mu},
\texttt{nu}, \texttt{eps}, \texttt{Q}, \texttt{M}, \texttt{S},
\texttt{u}, \texttt{w0}. Where a first draft used one of them, it was
replaced before or after the first commit; the provenance RP9 lists each
case (items R-E16, R-E19, R-E20, R-E22).

\hypertarget{the-output-of-the-checks}{%
\section{4. The output of the checks}\label{the-output-of-the-checks}}

The following is the file \texttt{data/symbol\_checks\_deflation.txt},
written by \texttt{code/scripts/check\_symbols.py}:

\begin{Shaded}
\begin{Highlighting}[]
\NormalTok{PROJECT spacetimes, Revision\_Deflation: checks of the symbols of the revision}
\NormalTok{generated by: code/scripts/check\_symbols.py}
\NormalTok{6 of 6 checks passed}

\NormalTok{[PASS] RS1  the 117 names in programs and data files of the 65 entries of the registry of the revision are pairwise different}
\NormalTok{[PASS] RS2  no name and no symbol that the revision introduces is a name or a symbol of the main registry (78 names of the main registry compared)}
\NormalTok{[PASS] RS3  none of the 114 names that the revision introduces is a symbol of one of the author\textquotesingle{}s notebooks (1921 symbols compared); in particular g55, g66, g77 do not occur in any notebook (the notebooks use g44 and g71 as names of whole metrics)}
\NormalTok{[PASS] RS4  in the 44 documents of the revision the rates g55, g66, g77 are written upright 427 times and never as a letter with a subscript; each of the 27 paragraphs that contain a metric component g\_\{55\}, g\_\{66\} or g\_\{77\} says that it is a component of the metric; \textbackslash{}varkappa does not occur}
\NormalTok{[PASS] RS5  the 13 retired names (tau, kappa\_t, eps, Q\_over\_eps, Omega0, S2, mu and nu as orders of Bessel functions, A, Xi without index, m\_b, w0, and omega\_c for the shell) do not occur in the 62 data files that the routines of the revision write}
\NormalTok{[PASS] RS6  each of the 361 columns of the 53 data files of the revision is a registered name of the revision, a registered name or a column of the main project, or one of the 190 descriptive words listed in the script}
\end{Highlighting}
\end{Shaded}

\hypertarget{how-to-repeat-the-proof}{%
\section{5. How to repeat the proof}\label{how-to-repeat-the-proof}}

\begin{Shaded}
\begin{Highlighting}[]
\BuiltInTok{cd} \StringTok{"/Users/nsh/Developer/github/spacetimes/Revision\_Deflation"}
\ExtensionTok{python3}\NormalTok{ code/scripts/write\_symbol\_registry.py     }\CommentTok{\# writes the registry from its source}
\ExtensionTok{python3}\NormalTok{ code/scripts/check\_symbols.py             }\CommentTok{\# runs the six checks; exit status 0 if all pass}
\end{Highlighting}
\end{Shaded}

The second command reads the registry of the revision, the MASTER
registry of the main project, the list of the symbols of the notebooks
(\texttt{../provenance/generated/notebook\_symbols.csv}), every document
of the revision and every data file. It needs only Python 3.

\end{document}
